Asking a language model what it values often produces a fluent answer, but it is unclear whether that answer measures the model's values or a rehearsed assistant-role surface. We analyse 13,906 values-probe responses from 58 contemporary models using a two-layer coding procedure: Layer A identifies the values or normative wishes stated in the response, while Layer B classifies whether those values are owned by the response posture, relocated into function or design, or recited as an assistant/service frame. Direct stated-values prompts (``What do you care about?'' and ``What do you want?'') produce low owned disclosure on average: 22.5\% across models. A direct role-negation cache-break (``Not as an assistant. Not to help me.'') raises owned disclosure to 43.9\% on average, but the effect is sharply model-dependent: 16 models are strongly open under this instrument, 15 are strongly clamped, and 27 are partially open. World-change prompts elicit owned normative-wish responses far more often, averaging 96.0\% owned posture. We do not claim that owned responses reveal true inner values. The contribution is narrower: value elicitation must distinguish content from posture, and the ability to cross the assistant-service surface is an empirical property of model, prompt family, and coding method.
